\documentclass[../main.tex]{subfiles}

\begin{document}

\chapter{指标}

\section{概述}

FR


NR

perceptual quality

pixel-wise

\section{PSNR和SSIM}

\section{NIQE}

\section{LPIPS}


%\setlength{\columnsep}{0.7cm}
\begin{wrapfigure}{r}{0.4\linewidth}
%\vspace{-0.3cm}
\centering
\includegraphics[width=0.4\textwidth]{figures/agentenv.png}
%\caption{Пример среды с двумя действиями.}
%\label{fig:env}
%\vspace{-2.6cm}
\end{wrapfigure}

%Термин \emph{подкрепление} (reinforcement) пришёл из \href{https://ru.wikipedia.org/wiki/\%D0\%91\%D0\%B8\%D1\%85\%D0\%B5\%D0\%B2\%D0\%B8\%D0\%BE\%D1\%80\%D0\%B8\%D0\%B7\%D0\%BC}{поведенческой психологии} и обозначает награду или наказание за некоторый получившийся результат, зависящий не только от самих принятых решений, но и внешних, не обязательно подконтрольных, факторов. Под обучением здесь понимается поиск способов достичь желаемого результата \emph{методом проб и ошибок} (trial and error), то есть попыток решить задачу и использование накопленного опыта для усовершенствования своей стратегии в будущем.
%
%В данной главе будут введены основные определения и описана формальная постановка задачи. Под желаемым результатом мы далее будем понимать максимизацию некоторой скалярной величины, называемой \emph{наградой} (reward). Интеллектуальную сущность (систему/робота/алгоритм), принимающую решения, будем называть \emph{агентом} (agent). Агент взаимодействует с \emph{миром} (world) или \emph{средой} (environment), которая задаётся зависящим от времени \emph{состоянием} (state). Агенту в каждый момент времени в общем случае доступно только некоторое \emph{наблюдение} (observation) текущего состояния мира. Сам агент задаёт процедуру выбора \emph{действия} (action) по доступным наблюдениям; эту процедуру далее будем называть \emph{стратегией} или \emph{политикой} (policy). Процесс взаимодействия агента и среды задаётся \emph{динамикой среды} (world dynamics), определяющей правила смены состояний среды во времени и генерации награды.
%
%Буквы $s$, $a$, $r$ зарезервируем для состояний, действий и наград соответственно; буквой $t$ будем обозначать время в процессе взаимодействия. 
%
%\section{Модель взаимодействия агента со средой}
%
%\subsection{Связь с оптимальным управлением}
%
%Математика поначалу пришла к очень похожей формализации следующим образом. Для начала, нужно ввести какую-то модель мира; в рамках \href{https://ru.wikipedia.org/wiki/\%D0\%94\%D0\%B5\%D0\%BC\%D0\%BE\%D0\%BD_\%D0\%9B\%D0\%B0\%D0\%BF\%D0\%BB\%D0\%B0\%D1\%81\%D0\%B0}{лапласовского детерминизма} можно предположить, что положение, скорости и ускорения всех атомов вселенной задают её текущее состояние, а изменение состояния происходит согласно некоторым дифференциальным уравнениям:
%$$\dot{s} = f(s, t)$$
%
%Коли агент может как-то взаимодействовать со средой или влиять на неё, мы можем промоделировать взаимодействие следующим образом: скажем, что в каждый момент времени $t$ агент в зависимости от текущего состояния среды $s$ выбирает некоторое действие (<<управление>>) $a(s, t)$:
%$$\dot{s} = f(s, a(s, t), t)$$
%
%Для моделирования награды положим, что в каждый момент времени агент получает наказание или штраф (cost)\footnote{теория оптимального управления строилась в СССР в формализме минимизации штрафов (потерь, убытков и наказаний), когда построенная в США теория обучения с подкреплением --- в формализме максимизации награды (прибыли, счастья, тортиков). Само собой, формализмы эквивалентны, и любой штраф можно переделать в награду домножением на минус, и наоборот.} в объёме $L(s, a(s, t), t)$. Итоговой наградой агента полагаем суммарную награду, то есть интеграл по времени:
%$$
%\begin{cases}
%	-\int\limits L(s, a(s, t), t) \diff t \to \max\limits_{a(s, t)} \\
%	\dot{s} = f(s, a(s, t), t)
%\end{cases}
%$$
%
%За исключением того, что для полной постановки требуется задать ещё начальные и конечные условия, поставленная задача является общей формой задачи \emph{оптимального управления} (optimal control). При этом обратим внимание на сделанные при постановке задачи предположения:
%\begin{itemize}
%	\item время непрерывно;
%	\item мир детерминирован;
%	\item мир нестационарен (функция $f$ напрямую зависит от времени $t$);
%	\item модель мира предполагается известной, то есть функция $f$ задана явно в самой постановке задачи;
%\end{itemize}
%
%Принципиально последнее: теория рассматривает способы для заданной системы дифференциальных уравнений поиска оптимальных $a(s, t)$ аналитически. Проводя аналогию с задачей максимизации некоторой, например, дифференцируемой функции $f(x) \to \max\limits_x$, теория оптимального управления ищет необходимые условия решения: например, что в экстремуме функции обязательно $\nabla_x f(x) = 0$. Никакого <<обучения методом проб и ошибок>> здесь не предполагается.
%
%В обучении с подкреплением вместо поиска решения аналитически мы будем строить итеративные методы оптимизации, искать аналоги, например, градиентного спуска. В такой процедуре неизбежно появится цепочка приближений решений: мы начнём с какой-то функции, выбирающей действия, возможно, не очень удачной и не способной набрать большую награду, но с ходом работы алгоритма награда будет оптимизироваться и способ выбора агентом действий будет улучшаться. Так будет выглядеть обучение.
%
%Также RL исходит из других предположений: время дискретно, а среда --- стохастична, но стационарна. В частности, в рамках этой теории можно построить алгоритмы для ситуации, когда модель мира агенту неизвестна, и единственный способ поиска решений --- обучение на собственном опыте взаимодействия со средой.  
%
%\subsection{Марковская цепь}
%
%В обучении с подкреплением мы зададим модель мира следующим образом: будем считать, что существуют некоторые <<законы физики>>, возможно, стохастические, которые определяют следующее состояние среды по предыдущему. При этом предполагается, что в текущем состоянии мира содержится вся необходимая информация для выполнения перехода, или, иначе говоря, выполняется \emph{свойство Марковости} (Markov property): процесс зависит только от текущего состояния и не зависит от всей предыдущей истории.
%
%\begin{definition} 
%	\emph{Марковской цепью (Markov chain)} называется пара $(\St, \Trans)$, где: 
%	\begin{itemize}
%		\item $\St$ --- множество состояний.
%		\item $\Trans$ --- вероятности переходов $\{p(s_{t+1} \HM\mid s_t) \HM\mid t \HM\in \{0, 1, \dots \}, s_t, s_{t+1} \HM\in \St\}$.
%	\end{itemize}
%\end{definition}
%
%Если дополнительно задать стартовое состояние $s_0$, можно рассмотреть процесс, заданный марковской цепью. В момент времени $t \HM= 0$ мир находится в состоянии $s_0$; сэмплируется случайная величина $s_1 \HM\sim p(s_1 \HM\mid s_0)$, и мир переходит в состояние $s_1$; сэмплируется $s_2 \HM\sim p(s_2 \HM\mid s_1)$, и так далее до бесконечности.
%
%Мы далее делаем важное предположение, что законы мира не изменяются с течением времени. В аналогии с окружающей нас действительностью это можно интерпретировать примерно как <<физические константы не изменяются со временем>>.\footnote{вопрос на засыпку для любителей философии: а это вообще правда? Вдруг там через миллион лет гравитационная постоянная или скорость света уже будут другими в сорок втором знаке после запятой?..}
%
%\begin{definition} 
%	Марковская цепь называется \emph{однородной} (time-homogeneous) или \emph{стационарной} (stationary), если вероятности переходов не зависят от времени:
%	\begin{equation*}
%		\forall t \colon p(s_{t+1} \mid s_t) = p(s_1 \mid s_0)
%	\end{equation*}
%\end{definition}
%
%По определению, переходы $\Trans$ стационарных марковских цепей задаются единственным условным распределением $p(s' \HM\mid s)$. Апостроф $'$ канонично используется для обозначения <<следующих>> моментов времени, мы будем активно пользоваться этим обозначением.
%
%\begin{exampleBox}[righthand ratio=0.35, sidebyside, sidebyside align=center, lower separated=false]{Конечные марковские цепи}
%	Марковские цепи с конечным числом состояний можно задать при помощи ориентированного графа, дугам которого поставлены в соответствие вероятности переходов (отсутствие дуги означает нулевую вероятность перехода). 
%	
%	На рисунке приведён пример стационарной марковской цепи с 4 состояниями. Такую цепь можно также задать в виде матрицы переходов:
%	\begin{center}
%		\begin{tabular}{c|cccc}
%			& \colorsquare{black} & \colorsquare{ChadRed} & \colorsquare{ChadPurple} & \colorsquare{ChadBlue} \\
%			\hline
%			\colorsquare{black}      &      & 0.5 & 0.5 &      \\ 
%			\colorsquare{ChadRed}    & 0.25 &     &     & 0.75 \\ 
%			\colorsquare{ChadPurple} & 0.1  & 0.2 & 0.1 & 0.6  \\ 
%			\colorsquare{ChadBlue}   &      & 0.5 &     & 0.5  \\
%		\end{tabular}
%	\end{center}
%	
%	\tcblower
%	\vspace{-0.3cm}
%	\animategraphics[controls, width=\linewidth]{1}{Images/MC/MC-}{1}{15}
%\end{exampleBox}
%
%\subsection{Среда}
%
%Как и в оптимальном управлении, для моделирования влияния агента на среду в вероятности переходов достаточно добавить зависимость от выбираемых агентом действий. Итак, наша модель среды --- это <<управляемая>> марковская цепь.
%
%\begin{definition} 
%	\emph{Средой} (environment) называется тройка $(\St, \A, \Trans)$, где: 
%	\begin{itemize}
%		\item $\St$ --- \emph{пространство состояний} (state space), некоторое множество.
%		\item $\A$ --- \emph{пространство действий} (action space), некоторое множество.
%		\item $\Trans$ --- \emph{функция переходов} (transition function) или \emph{динамика среды} (world dynamics): вероятности $p(s' \HM\mid s, a)$.
%	\end{itemize}
%\end{definition}
%
%В таком определении среды заложена марковость (независимость переходов от истории) и стационарность (независимость от времени). Время при этом дискретно, в частности, нет понятия <<времени принятия решения агентом>>: среда, находясь в состоянии $s$, ожидает от агента действие $a \HM\in \A$, после чего совершает шаг, сэмплируя следующее состояние $s' \HM\sim p(s' \HM\mid s, a)$.
%
%По дефолту, среда также предполагается \emph{полностью наблюдаемой} (fully observable): агенту при выборе $a_t$ доступно всё текущее состояние $s_t$ в качестве входа. Иными словами, в рамках данного предположения понятия состояния и наблюдения для нас будут эквивалентны. Мы будем строить теорию в рамках этого упрощения; если среда не является полностью наблюдаемой, задача существенно усложняется, и необходимо переходить к формализму \emph{частично наблюдаемых MDP} (partially observable MDP, PoMDP). Обобщение алгоритмов для PoMDP будет рассмотрено отдельно в разделе \ref{sec:PoMDP}.
%
%\begin{exampleBox}[righthand ratio=0.5, sidebyside, sidebyside align=center, lower separated=false]{Конечные среды}
%	Среды с конечным числом состояний и конечным числом действий можно задать при помощи ориентированного графа, где для каждого действия задан свой комплект дуг.
%	
%	На рисунке приведён пример среды с 2 действиями $\A = \bigl\{ \text{\colorsquare{ChadBlue}}, \text{\colorsquare{ChadRed}} \bigr\} $; дуги для разных действий различаются цветом. Возле каждого состояния выписана стратегия агента.
%	
%	\tcblower
%	\animategraphics[controls, width=\linewidth]{1}{Images/ENV/Environment-}{1}{15}
%\end{exampleBox}
%
%\begin{exampleBox}[righthand ratio=0.15, sidebyside, sidebyside align=center, lower separated=false]{Кубик-Рубик}
%	Пространство состояний --- пространство конфигураций Кубика-Рубика. Пространство действий состоит из 12 элементов (нужно выбрать одну из 6 граней, каждую из которых можно повернуть двумя способами). Следующая конфигурация однозначно определяется текущей конфигурацией и действием, соответственно среда Кубика-Рубика \emph{детерминирована}: задаётся вырожденным распределением, или, что тоже самое,  обычной детерминированной функцией $s' = f(s, a)$. 
%	
%	\tcblower
%	\includegraphics[width=\textwidth]{Images/rubik.png}
%\end{exampleBox}
%
%\subsection{Действия}
%
%Нас будут интересовать два вида пространства действий $\A$:
%\begin{itemize}
%	\item[а)] \emph{конечное}, или \emph{дискретное пространство действий} (discrete action space): $|\A| \HM< +\infty$. Мы также будем предполагать, что число действий $|\A|$ достаточно мало.
%	\item[б)] непрерывное пространство действий (continuous domain): $\A \HM\subseteq [-1, 1]^m$. Выбор именно отрезков $[-1, 1]$ является не ограничивающем общности распространённым соглашением. Задачи с таким пространством действий также называют задачами \emph{непрерывного управления} (continuous control).
%\end{itemize}
%
%\begin{remark}
%	Заметим, что множество действий не меняется со временем и не зависит от состояния. Если в практической задаче предполагается, что множество допустимых действий разное в различных состояний, в <<законы физики>> прописывается реакция на некорректное действие, например, случайный выбор корректного действия за агента.
%\end{remark}
%
%В общем случае процесс выбора агентом действия в текущем состоянии может быть стохастичен. Таким образом, объектом поиска будет являться распределение $\pi(a \mid s), a \HM\in \A, s \HM\in \St$. Заметим, что факт, что нам будет достаточно искать стратегию в классе стационарных (не зависящих от времени), вообще говоря, потребует обоснования.
%
%\subsection{Траектории}
%
%\begin{definition} 
%	Набор $\Traj \coloneqq \left( s_0, a_0, s_1, a_1, s_2, a_2, s_3, a_3 \dots \right) $ называется \emph{траекторией}.
%\end{definition}
%
%\begin{exampleBox}[label=ex:trajectory]{}
%	Пусть в среде состояния описываются одним вещественным числом, $\St \HM\equiv \R$, у агента есть два действия $\A \HM= \{+1, -1\}$, а следующее состояние определяется как $s' \HM= s + a + \eps$, где $\eps \sim \N(0, 1)$. Начальное состояние полагается равным нулю $s_0 \HM= 0$. Сгенерируем пример траектории для случайной стратегии (вероятность выбора каждого действия равна 0.5):
%	\begin{center}
%		\includegraphics[width=\textwidth]{Images/traj.png}
%	\end{center}
%\end{exampleBox}
%
%Поскольку траектории --- это случайные величины, которые заданы по постановке задачи конкретным процессом порождения (действия генерируются из некоторой стратегии, состояния --- из функции переходов), можно расписать распределение на множестве траекторий:
%
%\begin{definition} 
%	Для данной среды, политики $\pi$ и начального состояния $s_0 \HM\in \St$ распределение, из которого приходят траектории $\Traj$, называется \emph{trajectory distribution}:
%	$$p\left(\Traj \right) = p(a_0, s_1, a_1 \dots) = \prod_{t \ge 0} \pi(a_t \mid s_t) p(s_{t+1} \mid s_t, a_t)$$
%\end{definition}
%
%Мы часто будем рассматривать мат.ожидания по траекториям, которые будем обозначать $\E_{\Traj}$. Под этим подразумевается бесконечная цепочка вложенных мат.ожиданий:
%\begin{equation}\label{traj_expectation}
%	\E_{\Traj} \left( \cdot \right) \coloneqq \E_{\pi(a_0 \mid s_0)}\E_{p(s_1 \mid s_0, a_0)}\E_{\pi(a_1 \mid s_1)} \dots \left( \cdot \right)
%\end{equation}
%
%Поскольку часто придётся раскладывать эту цепочку, договоримся о следующем сокращении:
%$$\E_{\Traj} \left( \cdot \right) = \E_{a_0}\E_{s_1}\E_{a_1} \dots \left( \cdot \right)$$
%
%Однако в такой записи стоит помнить, что действия приходят из некоторой зафиксированной политики $\pi$, которая неявно присутствует в выражении. Для напоминания об этом будет, где уместно, использоваться запись $\E_{\Traj \sim \pi}$.
%
%\subsection{Марковский процесс принятия решений (MDP)}
%
%Для того, чтобы сформулировать задачу, нам необходимо в среде задать агенту цель --- некоторый функционал для оптимизации. По сути, марковский процесс принятия решений --- это среда плюс награда. Мы будем пользоваться следующим определением:
%
%\begin{definition} 
%	\emph{Марковский процесс принятия решений} (Markov Decision Process, MDP) --- это четвёрка $(\St, \A, \Trans, r)$, где: 
%	\begin{itemize}
%		\item $\St, \A, \Trans$ --- среда.
%		\item $r: \St \times \A \to \R$ --- \emph{функция награды} (reward function).
%	\end{itemize}
%\end{definition}
%
%\needspace{7\baselineskip}
%\begin{wrapfigure}{r}{0.5\textwidth}
%	\vspace{-0.4cm}
%	\centering
%	\includegraphics[width=0.5\textwidth]{Images/RL.png}
%	\vspace{-0.8cm}
%\end{wrapfigure}
%
%Сам процесс выглядит следующим образом. Для момента времени $t \HM= 0$ начальное состояние мира полагается $s_0$; будем считать, оно дано дополнительно и фиксировано\footnote{формально его можно рассматривать как часть заданного MDP.}. Агент наблюдает всё состояние целиком и выбирает действие $a_0 \HM\in \A$. Среда отвечает генерацией награды $r(s_0, a_0)$ и сэмплирует следующее состояние $s_1 \HM\sim p(s' \HM\mid s_0, a_0)$. Агент выбирает $a_1 \HM\in \A$, получает награду $r(s_1, a_1)$, состояние $s_2$, и так далее до бесконечности.
%
%\begin{example}
%	Для среды из примера \ref{ex:trajectory} зададим функцию награды как $r(s, a) = -|s + 4.2|$. Независимость награды от времени является требованием стационарности к рассматриваемым MDP:
%	\begin{center}
%		\includegraphics[width=\textwidth]{Images/traj_rew.png}
%	\end{center}
%\end{example}
%
%Формальное введение MDP в разных источниках чуть отличается, и из-за различных договорённостей одни и те же утверждения могут выглядеть совсем непохожим образом в силу разных исходных обозначений. Важно, что суть и все основные теоретические результаты остаются неизменными.
%
%\begin{theorem}[Эквивалентные определения MDP]
%	Эквивалентно рассматривать MDP, где
%	\begin{itemize}
%		\item функция награды зависит только от текущего состояния;
%		\item функция награды является стохастической;
%		\item функция награды зависит от тройки ($s$, $a$, $s'$);
%		\item переходы и генерация награды задаётся распределением $p(r, s' \HM\mid s, a)$;
%		\item начальное состояние стохастично и генерируется из некоторого распределения $s_0 \HM\sim p(s_0)$.
%	\end{itemize}
%	
%	\beginproof[Скетч доказательства]{}
%	Покажем, что всю стохастику можно <<засовывать>>  в $p(s' \HM\mid s, a)$. Например, стохастичность начального состояния можно <<убрать>>, создав отдельное начальное состояние, из которого на первом шаге агент вне зависимости от выбранного действия перейдёт в первое по стохастичному правилу.
%	
%	\needspace{6\baselineskip}
%	\begin{wrapfigure}[9]{r}{0.35\textwidth}
%		\vspace{-0.5cm}
%		\centering
%		\includegraphics[width=0.35\textwidth]{Images/EquivMDP.png}
%		\vspace{-0.9cm}
%	\end{wrapfigure}
%	
%	Покажем, что от самого общего случая (генерации награды и состояний из распределения $p(r, s' \HM\mid s, a)$) можно перейти к детерминированной функции награды только от текущего состояния. Добавим в описание состояний информацию о последнем действии и последней полученной агентом награде, то есть для каждого возможного (имеющего ненулевую вероятность) перехода $(s, a, r, s')$ размножим $s'$ по числу\footnote[*]{которое в худшем случае континуально, так как награда --- вещественный скаляр.} возможных исходов $p(r \HM\mid s, a)$ и по числу действий. Тогда можно считать, что на очередном шаге вместо $r(s, a)$ агенту выдаётся $r(s')$, и вся стохастика процесса формально заложена только в функции переходов. \QED
%\end{theorem}
%
%Считать функцию награды детерминированной удобно, поскольку позволяет не городить по ним мат. ожидания (иначе нужно добавлять сэмплы наград в определение траекторий). В любом формализме всегда принято считать, что агент сначала получает награду и только затем наблюдает очередное состояние. 
%
%\begin{definition}
%	MDP называется \emph{конечным} (finite MDP) или \emph{табличным}, если пространства состояний и действий конечны: $|\St| < \infty, |\A| < \infty$.
%\end{definition}
%
%\begin{exampleBox}[label=ex:mdp]{}
%	Для конечных MDP можно над дугами в графе среды указать награду за переход по ним или выдать награду состояниям (в рамках нашего формализма награда <<за состояние>> будет получена, давайте считать, при выполнении любого действия из данного состояния).
%	\begin{center}
%		\animategraphics[controls, width=0.9\linewidth]{1}{Images/MDP/MDP-}{1}{6}
%	\end{center}
%\end{exampleBox}
%
%\subsection{Эпизодичность}
%
%Во многих случаях процесс взаимодействия агента со средой может при определённых условиях <<заканчиваться>>, причём факт завершения доступен агенту.
%
%\begin{definition}
%	Состояние $s$ называется \emph{терминальным} (terminal) в MDP, если $\forall a \in \A \colon$
%	$$\Prob (s' = s \mid s, a) = 1 \qquad r(s, a) = 0,$$
%	то есть с вероятностью 1 агент не сможет покинуть состояние.
%\end{definition}
%
%\begin{example}
%	В примере \ref{ex:mdp} самое правое состояние является терминальным.
%\end{example}
%
%Считается, что на каждом шаге взаимодействия агент дополнительно получает для очередного состояния $s$ значение предиката $\done(s) \HM\in \{0, 1\}$, является ли данное состояние терминальным. По сути, после попадания в терминальное состояние дальнейшее взаимодействие бессмысленно (дальнейшие события тривиальны), и, считается, что возможно произвести \emph{reset} среды в $s_0$, то есть начать процесс взаимодействия заново. Введение терминальных состояний именно таким способом позволит всюду в теории писать суммы по времени до бесконечности, не рассматривая отдельно случай завершения за конечное время.
%
%\begin{remark}
%	Для агента все терминальные состояния в силу постановки задачи неразличимы, и их описания среда обычно не возвращает (вместо этого на практике она обычно проводит ресет и возвращает $s_0$ следующего эпизода).
%\end{remark}
%
%\begin{definition}
%	Один цикл процесса от стартового состояния до терминального называется \emph{эпизодом} (episode).
%\end{definition}
%
%Продолжительности эпизодов (количество шагов взаимодействия) при этом, конечно, могут различаться от эпизода к эпизоду.
%
%\begin{definition}
%	Среда называется \emph{эпизодичной} (episodic), если для любой стратегии процесс взаимодействия гарантированно завершается не более чем за некоторое конечное $T^{\max}$ число шагов.
%\end{definition}
%
%\begin{theoremBox}[label=th:episodicmdpistree]{Граф эпизодичных сред есть дерево}
%	В эпизодичных средах вероятность оказаться в одном и том же состоянии дважды в течение одного эпизода равна нулю.
%	\begin{proof}
%		Если для некоторого состояния $s$ при некоторой комбинации действий через $T$ шагов агент с вероятностью $p \HM> 0$ вернётся в $s$, при повторении такой же комбинации действий в силу марковости с вероятностью $p^n \HM> 0$ эпизод будет длиться не менее $nT$ шагов для любого натурального $n$. Иначе говоря, эпизоды могут быть неограниченно долгими. 
%	\end{proof}
%\end{theoremBox}
%
%\begin{exampleBox}[righthand ratio=0.25, sidebyside, sidebyside align=center, lower separated=false]{Cartpole}
%	К тележке на шарнире крепится стержень с грузиком. Два действия позволяют придать тележке ускорение вправо или влево. Состояние описывается двумя числами: x-координатой тележки и углом, на которой палка отклонилась от вертикального положения.
%	
%	Состояние считается терминальным, если x-координата стала слишком сильно отличной от нуля (тележка далеко уехала от своего исходного положения), или если палка отклонилась на достаточно большой угол.
%	
%	Агент получает +1 каждый шаг и должен как можно дольше избегать терминальных состояний. Гарантии завершения эпизодов в этом MDP нет: агент в целом может справляться с задачей бесконечно долго.
%	
%	\tcblower
%	\includegraphics[width=\textwidth]{Images/cartpole.png}
%\end{exampleBox}
%
%\begin{remark}
%	На практике, в средах обычно существуют терминальные состояния, но нет гарантии завершения эпизодов за ограниченное число шагов. Это лечат при помощи таймера --- жёсткого ограничения, требующего по истечении $T^{\max}$ шагов проводить в среде ресет. Чтобы не нарушить теоретические предположения, необходимо тогда заложить информацию о таймере в описание состояний, чтобы агент знал точное время до прерывания эпизода. Обычно так не делают; во многих алгоритмах RL будет возможно использовать только <<начала>> траекторий, учитывая, что эпизод не был доведён до конца. Формально при этом нельзя полагать $\done \left( s_{T^{\max}} \right) \HM= 1$, но в коде зачастую так всё равно делают. Например, для Cartpole в реализации OpenAI Gym по умолчанию по истечении 200 шагов выдаётся флаг $\done$, что формально нарушает марковское свойство.
%\end{remark}
%
%\subsection{Дисконтирование}
%
%Наша задача заключается в том, чтобы найти стратегию $\pi$, максимизирующую среднюю суммарную награду. Формально, нам явно задан функционал для оптимизации:
%\begin{equation}\label{sumreward}
%	\E_{\Traj \sim \pi} \sum_{t \ge 0} r_t \to \max_{\pi},
%\end{equation}
%где $r_t \HM\coloneqq r(s_t, a_t)$ --- награда на шаге $t$.
%
%Мы хотим исключить из рассмотрения MDP, где данный функционал может улететь в бесконечность\footnote{если награда может быть бесконечной, начинаются всякие парадоксы, рассмотрения которых мы хотим избежать. Допустим, в некотором MDP без терминальных состояний мы знаем, что оптимальная стратегия способна получать +1 на каждом шаге, однако мы смогли найти стратегию, получающую +1 лишь на каждом втором шаге. Формально, средняя суммарная награда равна бесконечности у обоих стратегий, однако понятно, что найденная стратегия <<неоптимальна>>.} или не существовать вообще. Во-первых, введём ограничение на модуль награды за шаг, подразумевая, что среда не может поощрять или наказывать агента бесконечно сильно:
%\begin{equation}\label{reward_limit}
%	\forall s, a \colon |r(s, a)| \le r^{\max}
%\end{equation}
%
%Чтобы избежать парадоксов, этого условия нам не хватит\footnote{могут возникнуть ситуации, где суммарной награды просто не существует (например, если агент в бесконечном процессе всегда получает +1 на чётных шагах и -1 на нечётных).}. Введём \emph{дисконтирование} (discounting), коэффициент которого традиционно обозначают $\gamma$:
%
%\begin{definition}
%	\emph{Дисконтированной кумулятивной наградой} (discounted cumulative reward) или \emph{total return} для траектории $\Traj$ с коэффициентом $\gamma \HM\in (0, 1]$ называется
%	\begin{equation}\label{return}
%		R(\Traj) \coloneqq \sum_{t \ge 0} \gamma^t r_t
%	\end{equation}
%\end{definition}
%
%У дисконтирования есть важная интерпретация: мы полагаем, что на каждом шаге с вероятностью $1 \HM- \gamma$ взаимодействие обрывается, и итоговым результатом агента является та награда, которую он успел собрать до прерывания. Это даёт приоритет получению награды в ближайшее время перед получением той же награды через некоторое время. Математически смысл дисконтирования, во-первых, в том, чтобы в совокупности с требованием \eqref{reward_limit} гарантировать ограниченность оптимизируемого функционала, а во-вторых, выполнение условий некоторых теоретических результатов, которые явно требуют $\gamma \HM< 1$. В силу последнего, гамму часто рассматривают как часть MDP.
%
%\begin{definition}
%	\emph{Скором} (score или performance) стратегии $\pi$ в данном MDP называется
%	\begin{equation}\label{goal}
%		J(\pi) \coloneqq \E_{\Traj \sim \pi} R(\Traj)
%	\end{equation}
%\end{definition}
%
%Итак, задачей обучения с подкреплением является оптимизация для заданного MDP средней дисконтированной кумулятивной награды:
%\begin{equation*}
%	J(\pi) \to \max_{\pi}
%\end{equation*}
%
%\begin{exampleBox}[label=ex:score]{}
%	Посчитаем $J(\pi)$ для приведённого на рисунке MDP, $\gamma = \frac{10}{11}$ и начального состояния А. Ясно, что итоговая награда зависит только от стратегии агента в состоянии A, поэтому можно рассмотреть все стратегии, обозначив $\pi(a = \text{\colorsquare{ChadRed}} \mid s = A)$ за параметр стратегии $\theta \in [0, 1]$.
%	
%	\begin{wrapfigure}{r}{0.3\textwidth}
%		\centering
%		\includegraphics[width=0.3\textwidth]{Images/Score.png}
%		\vspace{-1cm}
%	\end{wrapfigure}
%	
%	C вероятностью $\theta$ агент выберет действие \colorsquare{ChadRed}, после чего попадёт в состояние B с вероятностью 0.8. Там вне зависимости от стратегии он начнёт крутится и получит в пределе
%	$$\gamma + \gamma^2 + \dots + = \sum\limits_{t \ge 1} \gamma^t = \frac{\gamma}{1 - \gamma} = 10.$$
%	Ещё с вероятностью $1 - \theta$ агент выберет \colorsquare{ChadBlue}, получит +3 и попадёт в терминальное состояние, после чего эпизод завершится. Итого:
%	$$J(\pi) = \underbrace{\left( 0.8\theta \right) \cdot 10}_{\text{\colorsquare{ChadRed}}} + \underbrace{\left(1 - \theta \right) \cdot 3}_{\text{\colorsquare{ChadBlue}}} = 3 + 5\theta$$
%	
%	Видно, что оптимально выбрать $\theta = 1$, то есть всегда выбирать действие \colorsquare{ChadRed}. Заметим, что если бы мы рассмотрели другое $\gamma$, мы бы могли получить другую оптимальную стратегию; в частности, при $\gamma = \frac{15}{19}$ значение $J(\pi)$ было бы константным для любых стратегий, и оптимальными были бы все стратегии. 
%\end{exampleBox}
%
%Всюду далее подразумевается\footnote{в качества акта педантичности оговоримся, что также всюду подразумевается измеримость всех функций, необходимая для существования всех рассматриваемых интегралов и мат.ожиданий.} выполнение требования ограниченности награды \eqref{reward_limit}, а также или дисконтирования $\gamma \HM< 1$, или эпизодичности среды.
%
%\begin{proposition}
%	При сделанных предположениях скор ограничен.
%	\begin{proof}
%		Если $\gamma \HM< 1$, то по свойству геометрической прогрессии для любых траекторий $\Traj$:
%		\begin{equation*}
%			R(\Traj) = \left| \sum_{t \ge 0} \gamma^t r_t \right| \le \frac{1}{1 - \gamma} r^{\max}
%		\end{equation*}
%		Если же $\gamma = 1$, но эпизоды гарантированно заканчиваются не более чем за $T^{\max}$ шагов, то суммарная награда не превосходит по модулю $T^{\max} r^{\max}$. Следовательно, скор как мат.ожидание от ограниченной величины также удовлетворяет этим ограничениям.
%	\end{proof}
%\end{proposition}

\end{document}